\begin{helper}
Let \(C(x)=\binom{x}{2}_{\mathbb R}=x(x-1)/2\).  Suppose
\[
0\le\varepsilon\le1,\qquad 98\le D,\qquad 0\le b,\qquad 0\le u,\qquad
b+u=D.
\]
Then
\[
C((1+\varepsilon)b)+C((1+\varepsilon)u)
\le
(1+10\varepsilon)\bigl(C(b)+C(u)\bigr).
\]
\end{helper}
\begin{proof}
The identity
\[
2\Bigl((1+10\varepsilon)(C(b)+C(u))
-C((1+\varepsilon)b)-C((1+\varepsilon)u)\Bigr)
=
\varepsilon\Bigl((8-\varepsilon)(b^2+u^2)-9D\Bigr)
\]
follows by expanding \(C(x)=x(x-1)/2\) and using \(b+u=D\).  Since
\((b-u)^2\ge0\), the relation \(b+u=D\) gives
\[
D^2\le2(b^2+u^2).
\]
With \(D\ge98\) and \(\varepsilon\le1\), this implies
\[
(8-\varepsilon)(b^2+u^2)-9D\ge0.
\]
Multiplication by \(\varepsilon\ge0\) shows that the cleared denominator
difference is nonnegative, proving the claim.
\end{proof}
